\subsection{Педагогічні умови та методика навчання студентів розробки імерсивних освітніх ресурсів}\label{subsec:2.3}

Подальший хід дослідження зумовив необхідність визначення педагогічних умов, що створюють оптимальні передумови для формування у здобувачів умінь розробляти імерсивні освітні ресурси та інтегрувати їх у навчальне середовище.

Актуальність системного обґрунтування педагогічних умов підтверджується міжнародними дослідженнями: систематичний огляд К.~Ванг та Л.~Яньпін засвідчив, що відсутність теоретично обґрунтованого дизайну підготовки вчителів є ключовою слабкістю сучасних XR-програм~\cite{Wang2024XRTeacherEd}. Водночас Д.~Бек, Л.~Моргадо та П.~О'Ші, здійснивши мета-картування 47 оглядів імерсивного навчання, ідентифікували 45 стратегій і 21 практику, кластеризованих за напрямами: активний контекст, співпраця, залученість і scaffolding, присутність, мультимедійне навчання~\cite{Beck2024Practices}. Ці результати засвідчують необхідність створення системної рамки педагогічних умов, що інтегрує фактори-медіатори CAMIL~\cite{Makransky2021CAMIL} і компоненти TPACK у кожну умову.

Передусім визначимо сутність поняття <<педагогічні умови>> з огляду на вітчизняні педагогічні розвідки.

У довідкових джерелах поняття <<умова>> розуміється як необхідна обставина, що забезпечує можливість здійснення, створення або формування певного явища чи процесу та сприяє досягненню визначеного результату~\cite[c.~1506]{VTSSM2005}. Розвиваючи це загальнофілософське розуміння, дослідники конкретизують дефініцію у площині педагогічної науки: педагогічні умови трактуються як сприятливі педагогічні обставини та чинники, що свідомо створюються у закладі вищої освіти з метою цілеспрямованого впливу на здобувачів освіти~\cite[c.~32]{Yezhova2020Ref}. Водночас педагогічні умови осмислюються як узгоджена сукупність внутрішніх параметрів і зовнішніх характеристик організації та функціонування освітнього процесу, здатна забезпечити його високу результативність~\cite{Manko2000Conditions}. Концептуально близьку позицію висловлює Р.~Гуревич, пов'язуючи педагогічні умови з характеристиками процесу навчання та освітнього середовища~\cite{Gurevych2012Pedagogy}, а модернізація педагогічної освіти акцентує увагу на системному характері, структурованості та врахуванні особливостей готовності здобувачів до майбутньої професійної діяльності~\cite{Andrushchenko2018Prof}. Для студентів напряму підготовки <<Інформатика>> педагогічні умови розглядаються як обставини процесу навчання, в яких враховано наявні умови, передбачено способи їх перетворення та відібрано елементи змісту, методи й організаційні форми з урахуванням принципів оптимізації~\cite{Vdovychyn2015Conditions, Goncharenko2011Pedagogy, Pometun2004Interactive}.

Крім того, звертаємо увагу на класифікацію педагогічних умов (психолого-педагогічні, методичні, технологічні, організаційні), запропоновану в наукових джерелах, яку в межах методики навчання розробки імерсивних освітніх ресурсів майбутніх учителів інформатики адаптуємо відповідно до логіки та етапності цього процесу. Кожна з виокремлених умов співвіднесена з конкретними факторами-медіаторами моделі CAMIL~\cite{Makransky2021CAMIL} та компонентами TPACK, що забезпечує конструктивну узгодженість між теоретичними рамками (підрозділи~\ref{subsec:2.1}--\ref{subsec:2.2}) та практичною реалізацією. Таке узгодження адресує ключову прогалину, виявлену в XR-програмах підготовки вчителів~\cite{Wang2024XRTeacherEd}. Відповідність виокремлених педагогічних умов факторам CAMIL, компонентам TPACK і компонентам готовності (підрозділ~\ref{subsec:2.2}) подано на рис.~\ref{fig:conditions_alignment}.

\begin{figure}[!h]
\centering
\noindent\resizebox{\textwidth}{!}{%
\begin{tikzpicture}[every node/.style={font=\footnotesize}]

% Column headers
\node[dia/rbox=3.2cm, minimum height=1.0cm, font=\footnotesize\bfseries] (ch) {Педагогічна умова};
\node[dia/rbox=3.6cm, minimum height=1.0cm, font=\footnotesize\bfseries, right=0.15cm of ch] (cf) {Фактори CAMIL};
\node[dia/rbox=3.0cm, minimum height=1.0cm, font=\footnotesize\bfseries, right=0.15cm of cf] (ct) {Компоненти TPACK};
\node[dia/rbox=4.2cm, minimum height=1.0cm, font=\footnotesize\bfseries, right=0.15cm of ct] (cr) {Компоненти готовності};

% Row 1: Психолого-педагогічна
\node[dia/lrbox=3.2cm, below=0.15cm of ch] (r1c1) {психолого-педагогічна};
\node[dia/lrbox=3.6cm, below=0.15cm of cf] (r1c2) {мотивація, саморегуляція};
\node[dia/lrbox=3.0cm, below=0.15cm of ct] (r1c3) {PK, CK};
\node[dia/lrbox=4.2cm, below=0.15cm of cr] (r1c4) {мотиваційно-ціннісний, рефлексивно-оцінний};

% Row 2: Методична
\node[dia/lrbox=3.2cm, below=0.15cm of r1c1] (r2c1) {методична};
\node[dia/lrbox=3.6cm, below=0.15cm of r1c2] (r2c2) {когнітивне навантаження, агентність};
\node[dia/lrbox=3.0cm, below=0.15cm of r1c3] (r2c3) {PK, TPK, TPACK};
\node[dia/lrbox=4.2cm, below=0.15cm of r1c4] (r2c4) {когнітивний, діяльнісно-проєктувальний};

% Row 3: Технологічна
\node[dia/lrbox=3.2cm, below=0.15cm of r2c1] (r3c1) {технологічна};
\node[dia/lrbox=3.6cm, below=0.15cm of r2c2] (r3c2) {втілення, присутність};
\node[dia/lrbox=3.0cm, below=0.15cm of r2c3] (r3c3) {TK, TCK};
\node[dia/lrbox=4.2cm, below=0.15cm of r2c4] (r3c4) {технологічний, інтерактивний};

% Row 4: Організаційна
\node[dia/lrbox=3.2cm, below=0.15cm of r3c1] (r4c1) {організаційна};
\node[dia/lrbox=3.6cm, below=0.15cm of r3c2] (r4c2) {усі фактори (системна)};
\node[dia/lrbox=3.0cm, below=0.15cm of r3c3] (r4c3) {усі компоненти};
\node[dia/lrbox=4.2cm, below=0.15cm of r3c4] (r4c4) {усі компоненти};

\end{tikzpicture}%
}

\caption{Відповідність педагогічних умов факторам CAMIL, компонентам TPACK і компонентам готовності}
\label{fig:conditions_alignment}
\end{figure}

Підсумовуючи напрацювання вчених і з урахуванням особливостей предмета нашого дослідження, педагогічні умови підготовки майбутніх учителів інформатики до розробки ІОР доцільно конкретизувати крізь призму наведеної класифікації. У цьому контексті виокремлені нами умови корелюють передусім із психолого-педагогічною та методико-технологічною групами педагогічних умов і водночас відображають їх інтегративний характер.

Психолого-педагогічна умова~-- стимулювання професійної мотивації студентів через організацію навчального процесу, у межах якого створення та педагогічне застосування імерсивних освітніх ресурсів виступає засобом усвідомлення значущості майбутньої професійної діяльності, розвитку рефлексивного ставлення до професії та розкриття дидактичного потенціалу VR/AR-технологій у сучасній школі.

Методична умова~-- упровадження цілісної методики навчання, спрямованої на поетапне формування у студентів умінь проєктувати, педагогічно обґрунтовувати й інтегрувати імерсивні освітні ресурси в освітній процес відповідно до цілей і змісту шкільної інформатичної освіти.

Технологічна умова~-- забезпечення опанування студентами цифрових інструментів і середовищ для розробки імерсивних освітніх ресурсів, що уможливлює технічну реалізацію VR/AR-проєктів, їх тестування, удосконалення та адаптацію до освітніх потреб здобувачів загальної середньої освіти.

Організаційна умова~-- проєктування освітнього процесу на засадах поетапності, міждисциплінарної інтеграції та узгодженості цілей, змісту, методів і засобів навчання, а також визначення критеріїв, показників і рівнів сформованості готовності майбутніх учителів інформатики до розробки імерсивних освітніх ресурсів.

Проведений аналіз обумовив необхідність докладно охарактеризувати означені педагогічні умови, визначивши їхню значущість для ефективної організації підготовки майбутніх вчителів інформатики до створення імерсивних освітніх ресурсів.

\emph{\textbf{Психолого-педагогічна умова,}} що полягає в стимулюванні професійної мотивації студентів через організацію навчального процесу, де створення та педагогічне застосування імерсивних освітніх ресурсів стає потужним засобом усвідомлення значущості майбутньої професійної діяльності, набуває особливої ваги в підготовці майбутніх учителів інформатики. Вихідним положенням такого підходу виступає ідея, що практична діяльність з розробки VR/AR-ресурсів дозволяє студентам відчути реальний вплив цифрових технологій на освітній процес. Як засвідчує мета-огляд І.~Раду, доповнена реальність в освіті забезпечує зростання навчальних досягнень, підвищення мотивації та посилення залученості студентів, що є ключовими перевагами імерсивних технологій~\cite{KMotivation2023}. У контексті моделі CAMIL ця умова активізує передусім фактори \emph{мотивації} та \emph{саморегуляції}: самоефективність є найсильнішим предиктором прийняття VR-технологій майбутніми вчителями (PLS-SEM, \textit{N}\,=\,278)~\cite{Xie2024UTAUT2VR}, а безпосередній досвід створення імерсивних середовищ підвищує відчуття власної компетентності у сфері цифрової педагогіки. Коли майбутні вчителі самі створюють імерсивні середовища для уроків, вони не просто опановують інструменти~-- вони переживають, як ці інструменти змінюють сприйняття навчального матеріалу учнями, роблять навчання глибшим і емоційно насиченішим. Саме тому така діяльність уможливила зрушення в професійному самовизначенні студентів, адже вони бачать безпосередній зв'язок між власними зусиллями та потенційними успіхами школярів~\cite[c.~112]{Krokos2019Virtual}.

Провідну роль тут відіграє рефлексивне ставлення до професії. Під час створення імерсивних ресурсів студенти змушені постійно аналізувати власні дії: чому саме цей VR-сценарій ефективний, як він впливає на мотивацію учнів, які дидактичні принципи втілено найкраще. Як зазначають К.~Бошамп та К.~Томас, професійна ідентичність учителя формується саме через рефлексивне осмислення власної практики~\cite{ProfIdentity2024}. Мережевий епістемічний аналіз підтверджує, що дефіцити TPACK у планах уроків майбутніх учителів компенсуються саме через рефлексивний дискурс після проведення занять~\cite{Ga2024TPACKENA}, що валідує рефлексивний механізм, закладений у цій умові. Відповідно, створення імерсивних ресурсів стимулює цей процес: студенти починають сприймати себе не як користувачів технологій, а як творців інноваційного освітнього простору, що відповідає переходу від пасивного навчання до активного професійного зростання.

Організація навчального процесу навколо розробки та апробації імерсивних ресурсів забезпечило розкриття дидактичного потенціалу VR/AR-технологій у сучасній школі. Як засвідчує дослідження Дж.~Лі зі співавторами, саме діяльність з \emph{VR-Making} формує <<технологічну готовність, навички 4C та усвідомлення педагогічних переваг>> завдяки практичному створенню, а не лише споживанню контенту~\cite{Lee2022VRMaking}. Коли майбутні вчителі інформатики самі створюють такі ресурси, вони не лише вивчають теоретичні аспекти, а й переконуються на практиці в їхній ефективності. Це зумовило глибше розуміння того, як VR/AR може трансформувати традиційні уроки в інтерактивні, мотивуючі досвіди, що підтверджується дослідженням В.~Коваленко, М.~Мар'єнко та А.~Сухіх щодо застосування технологій доповненої та віртуальної реальності в умовах змішаного навчання~\cite{Pinchuk2022AR}. Такий підхід виявився особливо ефективним у формуванні стійкої професійної мотивації. Як доводять П.~Ертмер та А.~Оттенбрайт-Лефтвіч, готовність учителів до впровадження технологічних інновацій визначається сукупністю чинників~-- знань, самоефективності, педагогічних переконань і культури навчального закладу~\cite{KInnovation2023}. Відповідно, студенти, які набувають досвіду створення імерсивних проєктів, усвідомлюють, що володіння цими технологіями формує їхню конкурентоспроможність як фахівців, здатних відповідати викликам цифрової освіти.

Отже, організація навчального процесу, де створення та педагогічне застосування імерсивних освітніх ресурсів виступає ключовим елементом, уможливила не лише технічне оволодіння VR/AR, але й глибоке усвідомлення їхньої дидактичної цінності, розвиток рефлексії та стійкої мотивації до професійної діяльності. У термінах CAMIL ця умова забезпечує активізацію факторів \emph{мотивації} (через усвідомлення цінності технологій) та \emph{саморегуляції} (через рефлексивне оцінювання власних результатів), а в площині TPACK~-- розвиток PK (педагогічні переконання) та CK (розуміння предметного змісту).

\emph{\textbf{Методична умова,}} що полягає в упровадженні цілісної методики навчання, спрямованої на поетапне формування у здобувачів умінь проєктувати, педагогічно обґрунтовувати й інтегрувати імерсивні освітні ресурси в освітній процес відповідно до цілей і змісту шкільної інформатичної освіти, виявилася визначальною для якісної підготовки майбутніх учителів інформатики.

Вихідним положенням розробки методики став принцип системності формування професійних компетентностей. Спочатку здобувачі опановують базові технічні навички створення імерсивного контенту, потім переходять до аналізу дидактичних можливостей VR/AR, а на завершальному етапі вчаться інтегрувати розроблені ресурси в уроки інформатики, чітко узгоджуючи їх із програмними вимогами та віковими особливостями учнів. Така послідовність, на наше переконання, забезпечує уникнення фрагментарного засвоєння знань. Підхід <<Think-Create-Teach>>, верифікований на вибірці \textit{N}\,=\,108 студентів педагогічних спеціальностей, засвідчив, що дизайн-мислення як рамка для керованого створення навчальних матеріалів забезпечує суттєво вищу якість розроблених ресурсів порівняно з некерованим підходом~\cite{Calavia2023TCT}. Ця триетапна послідовність безпосередньо відповідає фазам 6-компонентної моделі проєктування ІОР (аналіз потреб → вибір ПЗ та розробка → впровадження та оцінювання)~\cite{Semerikov2022IERDesign}.

Визначальним елементом цієї методики постає педагогічна аргументація кожного розробленого імерсивного ресурсу. Здобувачі освіти опановують уміння критично осмислювати: яку дидактичну мету реалізує конкретний VR-сценарій, яким чином він забезпечує досягнення запланованих результатів навчання, які способи й прийоми стимулювання пізнавальної активності учнів у ньому закладені.

Цілісна методика уможливлює формування в здобувачів уміння в подальшому інтегрувати імерсивні ресурси саме в шкільну інформатичну освіту. У цьому контексті доцільним є звернення до п'яти перших принципів навчання, обґрунтованих М.~Меріллом: проблемно-центрованість, активізація попереднього досвіду, демонстрація, застосування та інтеграція~\cite{KGeneral2020}. Емпірична валідація цих принципів засвідчила значущо вищі навчальні досягнення та задоволеність у групі, де вони застосовувалися систематично (pre/post, \textit{N}\,=\,291)~\cite{Badali2020MerrillMOOC}. Реалізація цих принципів у процесі розробки ІОР забезпечує побудову навчального досвіду навколо автентичних проблем, поетапне залучення здобувачів до проєктної діяльності та інтеграцію набутих умінь у контекст майбутньої професійної практики. Практичне втілення методичної умови спирається на навчальний посібник із WebAR~\cite{Shepiliev2020WebAR} та методику інтеграції ML-моделей у WebAR-ресурси~\cite{Semerikov2024WebARML_CoSinE}, які забезпечують авторське науково-методичне підґрунтя. За таких умов імерсивні технології не просто <<прикрашають>> урок, а суттєво підвищують його ефективність і відповідність сучасним освітнім стандартам.

У термінах CAMIL методична умова керує фактором \emph{когнітивного навантаження} через scaffolding складності та розвиває \emph{агентність} студентів як активних дизайнерів освітніх ресурсів; у площині TPACK~-- забезпечує формування PK (педагогічна аргументація) та TPK (інтеграція технології з дидактикою).

Отже, упровадження цілісної поетапної методики навчання забезпечує не лише оволодіння здобувачами практичними навичками створення імерсивних ресурсів, але й сприяє сформуванню в них здатності свідомо, педагогічно обґрунтовано й системно інтегрувати ці технології в процес шкільної інформатичної освіти.

\emph{\textbf{Технологічна умова.}} Її дотримання спрямоване на формування у майбутніх учителів інформатики здатності комплексно організовувати процес технічного створення імерсивних освітніх продуктів, охоплюючи всі ключові етапи~-- від концептуалізації ідеї та цифрового проєктування до програмної реалізації, апробації, аналізу функціонування й подальшого вдосконалення розробленого ресурсу. Технологічна умова охоплює прогресивне опанування інструментів WebAR: MindAR (відстеження зображень), Three.js (3D-рендеринг), WebXR Device API (відстеження поверхонь), TensorFlow.js (ML-інтеграція)~\cite{Shepiliev2020WebAR, Semerikov2024WebARML_CoSinE}. У межах цієї умови акцент переноситься на роботу з різними імерсивними середовищами, типами інтерфейсів і способами взаємодії користувача з віртуальним контентом. Така підготовка дає змогу студентам усвідомлено добирати й модифікувати технологічні рішення з урахуванням вікових, когнітивних та освітніх особливостей здобувачів загальної середньої освіти.

Визначальним аспектом технологічної умови є модель прогресивної складності: від A-Frame (найпростіший рівень,~$\sim$10 рядків коду) через Three.js + MindAR (помірна складність) до TensorFlow.js + WebAR (просунутий рівень)~\cite{Semerikov2024WebAR_EduDim}. Така градація реалізує фактори CAMIL \emph{втілення} та \emph{присутності} через дедалі більш імерсивний досвід розробки: від візуалізації 3D-об'єктів до створення ML-інтегрованих інтерактивних середовищ. П'ять характеристик ІОР, визначених у підрозділі~\ref{subsec:1.2} (сенсорне занурення, інтерактивність, освітня спрямованість, адаптивність, вимірюваність), виступають технологічними орієнтирами на кожному рівні складності.

Водночас 15 доказових рекомендацій із проєктування AR-додатків, систематизованих за шістьма вимірами~\cite{GonnermannMuller2024ARDesign}, забезпечують конкретні технічні настанови для студентів. Міждисциплінарна співпраця між студентами педагогічних і технічних спеціальностей~\cite{Oberdorfer2021Interdisciplinary} гарантує одночасне забезпечення дидактичної та технічної якості розроблюваних ресурсів. Курс <<Синтетична реальність в освіті>> (18 модулів) реалізує цю модель: модулі 1--6 охоплюють основи WebAR та інструменти розробки, модулі 7--12~-- трекінг зображень, облич і поверхонь, модулі 13--18~-- ML-інтеграцію та проєктну діяльність.

Технологічна підготовка розглядається не як автономне опанування інструментів, а як складник цілісної професійної підготовки, інтегрований із методичними та педагогічними засадами навчання. У площині TPACK ця умова розвиває TK (технологічні знання) та TCK (технологічно-змістовну інтеграцію), а в термінах CAMIL~-- забезпечує фактори \emph{втілення} (тілесний досвід взаємодії з 3D-об'єктами) та \emph{присутності} (відчуття перебування в імерсивному середовищі). Така інтеграція забезпечує готовність майбутнього вчителя не лише використовувати VR/AR-рішення в освітньому процесі, а й гнучко трансформувати їх відповідно до зміни навчальних цілей, змісту теми або умов організації навчання.

\emph{\textbf{Організаційна умова}} передбачає цілеспрямоване проєктування освітнього процесу як логічно впорядкованої та внутрішньо узгодженої системи, у межах якої підготовка майбутніх учителів інформатики до розробки імерсивних освітніх ресурсів здійснюється поетапно, з урахуванням зростання складності навчальних завдань і поступового ускладнення професійних дій. Поетапність організації навчання забезпечує послідовний перехід від засвоєння базових уявлень про імерсивні технології та їх освітній потенціал до інтегрованого застосування знань і вмінь у процесі самостійного проєктування та педагогічного обґрунтування VR/AR-рішень. Цикл досвідного навчання Д.~Колба (конкретний досвід → рефлексивне спостереження → абстрактна концептуалізація → активне експериментування) забезпечує структурування практичних занять: лабораторна робота (CE) → аналіз результатів (RO) → теоретичне узагальнення (AC) → самостійний проєкт (AE)~\cite{Kee2023Experiential}.

Важливим складником цієї умови виступає міждисциплінарна інтеграція, що забезпечує взаємопроникнення змісту фахових, педагогічних і методичних дисциплін. Як підтвердили В.~Обердорфер зі співавторами, міждисциплінарний підхід у навчанні VR забезпечує <<взаємні вигоди>>: значущі приріст знань і в медіадидактиці, і в технології (pre/post)~\cite{Oberdorfer2021Interdisciplinary}. Такий підхід створює підґрунтя для формування цілісного бачення імерсивного освітнього ресурсу як результату поєднання технологічних, дидактичних і психолого-педагогічних компонентів. Унаслідок цього здобувачі освіти усвідомлюють взаємозалежність програмування, цифрового дизайну, методики навчання інформатики та вікових особливостей учнів, що сприяє подоланню фрагментарності підготовки й орієнтації виключно на технічний аспект.

Організаційна умова також передбачає узгодженість цілей, змісту, методів і засобів навчання, що унеможливлює стихійне або епізодичне використання імерсивних технологій у фаховій підготовці. Систематичний огляд XR-програм підготовки вчителів виявив, що <<відсутність конструктивної узгодженості між цілями, методами та оцінюванням>> є ендемічною слабкістю~\cite{Wang2024XRTeacherEd}~-- саме цю прогалину адресує організаційна умова. Навчальні цілі конкретизуються через очікувані результати, зміст добирається з урахуванням професійної значущості імерсивних ресурсів для шкільної інформатичної освіти, методи орієнтуються на активну, проєктну та рефлексивну діяльність студентів, а засоби навчання розглядаються як інструменти реалізації цілісного дидактичного задуму. Така узгодженість забезпечує логіку освітнього процесу та сприяє формуванню готовності до системного використання VR/AR-технологій у майбутній професійній діяльності.

Окрему роль у реалізації організаційної умови відіграє визначення чітких критеріїв, показників і рівнів сформованості готовності майбутніх учителів інформатики до розробки імерсивних освітніх ресурсів. Наявність діагностично обґрунтованого інструментарію дозволяє здійснювати моніторинг динаміки професійного становлення студентів, виявляти сильні сторони та проблемні зони підготовки, а також коригувати зміст і методи навчання відповідно до досягнутих результатів. У площині CAMIL організаційна умова є метаумовою, що інтегрує \emph{усі} фактори-медіатори системно, а в TPACK~-- забезпечує розвиток усіх компонентів. У такий спосіб організаційна умова забезпечує не лише структурну впорядкованість освітнього процесу, а й його результативність, орієнтовану на формування цілісної готовності до педагогічно вмотивованої та методично виваженої розробки імерсивних освітніх ресурсів.

Визначені педагогічні умови створюють рамку; методика навчання операціоналізує їх через конкретний триетапний процес, узгоджений із 6-компонентною моделлю проєктування ІОР~\cite{Semerikov2022IERDesign} та п'ятьма першими принципами навчання М.~Меріла~\cite{KGeneral2020}.

\emph{\textbf{Етап 1~-- теоретико-аналітичний}} відповідає фазам <<Аналіз потреб>> та <<Формулювання вимог>> 6-компонентної моделі. Здобувачі опановують теоретичні засади ІОР, п'ять їх сутнісних характеристик~\cite{Semerikov2022IERDesign}, здійснюють порівняльний аналіз інструментів WebAR (A-Frame, Three.js, MindAR, Babylon.js)~\cite{Shepiliev2020WebAR}. У термінах CAMIL цей етап керує фактором \emph{когнітивного навантаження} через scaffolding теоретичного матеріалу; у площині TPACK~-- закладає фундаментальні TK (знання інструментів) та CK (розуміння предметного змісту імерсивних технологій).

\emph{\textbf{Етап 2~-- практико-розробницький}} відповідає фазам <<Вибір програмного забезпечення>>, <<Планування>> та <<Розробка>>. Здобувачі реалізують прогресивну послідовність WebAR-проєктів: відстеження зображень → відстеження облич → відстеження поверхонь → ML-інтеграція~\cite{Semerikov2024WebARML_CoSinE}. Цикл Д.~Колба застосовується на кожному модулі: лабораторна робота (конкретний досвід) → аналіз результатів (рефлексивне спостереження) → теоретичне узагальнення (абстрактна концептуалізація) → самостійний проєкт (активне експериментування)~\cite{Kee2023Experiential}. У термінах CAMIL етап активізує \emph{втілення} (тактильна взаємодія з 3D-об'єктами), \emph{агентність} (студент як дизайнер) та \emph{присутність} (занурення у власні розробки); у площині TPACK~-- формує TPK (інтеграція технологій із дидактикою) та TCK (технологічно-змістовна інтеграція).

\emph{\textbf{Етап 3~-- інтегративно-рефлексивний}} відповідає фазам <<Впровадження>> та <<Оцінювання>>. Здобувачі проєктують цілісні ІОР-інтегровані уроки шкільної інформатики, здійснюють педагогічну аргументацію, проводять взаємне рецензування розроблених ресурсів. Квести профорієнтації із WebAR~\cite{Shepiliev2021Career} демонструють практичне застосування в реальних умовах. Оцінювання узгоджується з п'ятьма характеристиками ІОР (підрозділ~\ref{subsec:1.2}) та шістьма компонентами готовності (підрозділ~\ref{subsec:2.2}). У термінах CAMIL цей етап активізує \emph{саморегуляцію} через рефлексивне оцінювання; у площині TPACK~-- забезпечує повну інтеграцію TPACK як цілісного конструкту.

Поетапну послідовність методики, її відповідність фазам 6-компонентної моделі ІОР, факторам CAMIL, компонентам TPACK і модулям курсу <<Синтетична реальність в освіті>> подано на рис.~\ref{fig:methodology_stages}.

\begin{figure}[!h]
\centering
\noindent\resizebox{\textwidth}{!}{%
\begin{tikzpicture}[every node/.style={font=\footnotesize}]

% Three stage headers
\node[dia/header=4.5cm, minimum height=1.0cm] (s1h)
  {\small\textbf{Етап 1: Теоретико-аналітичний}};
\node[dia/header=4.5cm, minimum height=1.0cm, right=0.8cm of s1h] (s2h)
  {\small\textbf{Етап 2: Практико-розробницький}};
\node[dia/header=4.5cm, minimum height=1.0cm, right=0.8cm of s2h] (s3h)
  {\small\textbf{Етап 3: Інтегративно-рефлексивний}};

% Arrows between stages
\draw[dia/arrow] (s1h.east) -- (s2h.west);
\draw[dia/arrow] (s2h.east) -- (s3h.west);

% Stage 1 content boxes (placed first so labels can reference them)
\node[dia/lrbox=4.5cm, below=0.7cm of s1h.south, anchor=north] (s1r1) {%
  аналіз потреб,\\формулювання вимог};
\node[dia/lrbox=4.5cm, below=0.5cm of s1r1, anchor=north] (s1r2) {%
  когнітивне навантаження\\(scaffolding)};
\node[dia/lrbox=4.5cm, below=0.5cm of s1r2, anchor=north] (s1r3) {%
  TK + CK\\(базові знання)};
\node[dia/lrbox=4.5cm, below=0.5cm of s1r3, anchor=north] (s1r4) {%
  модулі 1--6\\(теорія WebAR, інструменти)};

% Stage 2 content boxes (aligned to Stage 1 rows)
\node[dia/lrbox=4.5cm, anchor=north] at (s2h.south |- s1r1.north) (s2r1) {%
  вибір ПЗ, планування,\\розробка};
\node[dia/lrbox=4.5cm, anchor=north] at (s2h.south |- s1r2.north) (s2r2) {%
  втілення, агентність,\\присутність};
\node[dia/lrbox=4.5cm, anchor=north] at (s2h.south |- s1r3.north) (s2r3) {%
  TPK, TCK\\(інтеграція)};
\node[dia/lrbox=4.5cm, anchor=north] at (s2h.south |- s1r4.north) (s2r4) {%
  модулі 7--12\\(трекінг, 3D, ML)};

% Stage 3 content boxes (aligned to Stage 1 rows)
\node[dia/lrbox=4.5cm, anchor=north] at (s3h.south |- s1r1.north) (s3r1) {%
  впровадження,\\оцінювання};
\node[dia/lrbox=4.5cm, anchor=north] at (s3h.south |- s1r2.north) (s3r2) {%
  саморегуляція,\\мотивація};
\node[dia/lrbox=4.5cm, anchor=north] at (s3h.south |- s1r3.north) (s3r3) {%
  TPACK\\(повна інтеграція)};
\node[dia/lrbox=4.5cm, anchor=north] at (s3h.south |- s1r4.north) (s3r4) {%
  модулі 13--18\\(проєкти, рефлексія)};

% Row labels on the left (positioned relative to Stage 1 content boxes)
\node[font=\footnotesize\bfseries\itshape, anchor=east] at ($(s1r1.west) + (-0.4cm,0)$) {Фази моделі ІОР:};
\node[font=\footnotesize\bfseries\itshape, anchor=east] at ($(s1r2.west) + (-0.4cm,0)$) {Фактори CAMIL:};
\node[font=\footnotesize\bfseries\itshape, anchor=east] at ($(s1r3.west) + (-0.4cm,0)$) {TPACK:};
\node[font=\footnotesize\bfseries\itshape, anchor=east] at ($(s1r4.west) + (-0.4cm,0)$) {Модулі курсу:};

% Dominant condition label at bottom
\node[dia/rbox=15.1cm, below=0.6cm of s2r4.south, font=\footnotesize] (cond) {%
  \textbf{Провідна умова:}~~психолого-педагогічна \textbar{} технологічна + методична \textbar{} організаційна%
};

\end{tikzpicture}%
}

\caption{Триетапна методика навчання розробки ІОР: відповідність фазам моделі, факторам CAMIL, компонентам TPACK і модулям курсу}
\label{fig:methodology_stages}
\end{figure}

Таким чином, триетапна методика навчання розробки імерсивних освітніх ресурсів (теоретико-аналітичний → практико-розробницький → інтегративно-рефлексивний) одночасно реалізує усі чотири педагогічні умови, операціоналізуючи їх через конкретний зміст, форми та засоби навчання. Курс <<Синтетична реальність в освіті>> (18 тижневих модулів) на базі навчального посібника з WebAR виступає конкретною реалізацією цієї методики, а її ефективність підлягає експериментальній перевірці (розділ~3).

